% ============================================================
%  利用 Warp Shuffle 实现高效并行归约 —— 以 Softmax 为例
%  张潇瀚 · 10245102417 · 并行计算课程项目
% ============================================================
\documentclass[12pt,a4paper]{ctexart}

% ---- Packages ----
\usepackage[top=2.5cm,bottom=2.5cm,left=2.8cm,right=2.8cm]{geometry}
\usepackage{amsmath,amssymb}
\usepackage{listings}
\usepackage{xcolor}
\usepackage{booktabs}
\usepackage{hyperref}
\usepackage{graphicx}
\usepackage{caption}
\usepackage{float}
\usepackage{enumitem}
\usepackage{fancyhdr}
\setlength{\headheight}{14.5pt}
\addtolength{\topmargin}{-2.5pt}

% ---- Font setup ----
% 等宽字体需支持制表符(├─│└)和 CJK 注释
\IfFontExistsTF{DejaVu Sans Mono}{
  \setmonofont{DejaVu Sans Mono}[Scale=MatchLowercase]       % Linux
}{\IfFontExistsTF{Consolas}{
  \setmonofont{Consolas}[Scale=MatchLowercase]               % Windows
}{\IfFontExistsTF{Menlo}{
  \setmonofont{Menlo}[Scale=MatchLowercase]                  % macOS
}{}}}
% CJK 等宽字体用于 listings 中的中文注释
\IfFontExistsTF{Noto Sans Mono CJK SC}{
  \setCJKmonofont{Noto Sans Mono CJK SC}[Scale=MatchLowercase]      % Linux
}{\IfFontExistsTF{NSimSun}{
  \setCJKmonofont{NSimSun}[Scale=1.05]                               % Windows 新宋体
}{\IfFontExistsTF{STHeiti Mono}{
  \setCJKmonofont{STHeiti Mono}[Scale=0.95]                          % macOS
}{}}}}

% ---- Code style ----
\definecolor{codebg}{RGB}{245,245,243}
\definecolor{codecomment}{RGB}{115,115,115}
\definecolor{codekeyword}{RGB}{0,47,167}
\definecolor{codefunc}{RGB}{196,37,107}
\definecolor{codestring}{RGB}{10,10,10}

\lstdefinestyle{cuda}{
    language=C++,
    basicstyle=\small,
    backgroundcolor=\color{codebg},
    commentstyle=\color{codecomment}\itshape,
    keywordstyle=\color{codekeyword}\bfseries,
    stringstyle=\color{codestring},
    numbers=left,
    numberstyle=\tiny\color{codecomment},
    frame=single,
    framerule=0.5pt,
    rulecolor=\color{codecomment!30},
    breaklines=true,
    breakatwhitespace=false,
    tabsize=4,
    showstringspaces=false,
    captionpos=b,
    aboveskip=10pt,
    belowskip=10pt,
    % CUDA kernel launch <<<...>>> confuses listings' C++ parser.
    % Use escapeinside to wrap them in safe mode:
    escapeinside={(*@}{@*)},
}
\lstdefinestyle{python}{
    language=Python,
    basicstyle=\small,
    backgroundcolor=\color{codebg},
    commentstyle=\color{codecomment}\itshape,
    keywordstyle=\color{codekeyword}\bfseries,
    stringstyle=\color{codestring},
    numbers=left,
    numberstyle=\tiny\color{codecomment},
    frame=single,
    framerule=0.5pt,
    breaklines=true,
    tabsize=4,
    showstringspaces=false,
    captionpos=b,
}
\lstdefinestyle{shell}{
    basicstyle=\footnotesize,
    backgroundcolor=\color{codebg},
    frame=single,
    framerule=0.5pt,
    breaklines=true,
    showstringspaces=false,
}

% ---- Headers ----
\pagestyle{fancy}
\fancyhf{}
\fancyhead[L]{利用 Warp Shuffle 实现高效并行归约}
\fancyhead[R]{张潇瀚 · 10245102417}
\fancyfoot[C]{\thepage}
\renewcommand{\headrulewidth}{0.4pt}

% ---- Metadata ----
\title{\textbf{利用 Warp Shuffle 实现高效并行归约} \\[8pt] \large —— 以 Softmax 为例}
\author{张潇瀚\ \ 10245102417}
\date{并行计算课程项目 \\[4pt] 2026年6月18日}

\begin{document}

% ============================================================
%  封面页（第 1 页）
% ============================================================
\begin{titlepage}
\thispagestyle{empty}
\begin{center}
\vspace*{4cm}

{\LARGE\textbf{利用 Warp Shuffle 实现高效并行归约}}\\[12pt]
{\Large —— 以 Softmax 为例}\\[3cm]

{\Large 张潇瀚}\\[6pt]
{\large 学号：10245102417}\\[2cm]

{\large 并行计算课程项目}\\[8pt]
{\large 2026年6月18日}\\[6cm]

\vfill
{\small GPU: NVIDIA GeForce RTX 4060 Laptop GPU $\cdot$ CUDA 13.3}
\end{center}
\end{titlepage}

% ============================================================
%  正文从第 2 页开始
% ============================================================
\newpage
\setcounter{page}{1}

\begin{abstract}
\noindent
本实验从零实现了两个 CUDA 版本的 Softmax 内核——串行 Naive 版本与基于 Warp Shuffle 的并行归约版本，封装为 PyTorch 自定义算子，并与 PyTorch 官方实现（底层 cuDNN）进行正确性交叉验证和性能对比。

实验结果表明：Warp Shuffle 归约在 $D \geq 256$ 时与 CUDA 自带的 cuDNN softmax 性能接近（差距 $<15\%$），$D=1024$ 时仅慢 $13\%$；所有测试偏差 $<2\times10^{-7}$；代码仅 80 行即达到接近硬件极限的性能。
\end{abstract}

\vspace{8pt}
\noindent\textbf{关键词：} CUDA；Warp Shuffle；并行归约；Softmax；GPU 编程

\newpage
\setcounter{page}{1}
\tableofcontents
\newpage

% ============================================================
\section{算法背景}
% ============================================================

\subsection{Softmax 定义}

Softmax 是神经网络中最常用的激活函数之一，将任意实数向量转换为概率分布：

\begin{equation}
\text{softmax}(x_i) = \frac{e^{x_i}}{\sum_{j} e^{x_j}}
\end{equation}

为避免 $e^{x}$ 溢出（float 最大值约 $3.4 \times 10^{38}$），实际采用 Safe Softmax：先找最大值 $m = \max_j x_j$，再计算：

\begin{equation}
\text{softmax}(x_i) = \frac{e^{x_i - m}}{\sum_j e^{x_j - m}}
\end{equation}

减去最大值后，最大项变为 $e^0 = 1$，永远不会溢出。

\subsection{并行归约}

Softmax 的核心操作包含\emph{两次归约（Reduction）}：

\begin{enumerate}[itemsep=4pt]
    \item \textbf{Max Reduction}：找 $D$ 个元素的最大值
    \item \textbf{Sum Reduction}：求 $D$ 个 $\exp$ 值的总和
\end{enumerate}

归约是并行计算中最基础的操作——把多个值合并为一个值。对于一个长度为 $D$ 的向量：
\begin{itemize}
    \item 串行归约需要 $O(D)$ 时间
    \item 用 $P$ 个处理器并行归约，理想情况下只需 $O(D/P + \log P)$ 时间
\end{itemize}

本实验的核心即利用 CUDA Warp Shuffle 指令在 GPU 上高效实现这两次归约。

% ============================================================
\section{系统架构}
% ============================================================

整套系统分为四层，调用链如下：

\begin{center}
\begin{tabular}{c}
\textbf{Python 用户接口} \\
\verb|from softmax import softmax_warp| \\
$\downarrow$ \\[4pt]
\textbf{C++ Torch 绑定} (\verb|softmax_bindings.cpp|, g++ 编译) \\
\verb|softmax_forward_warp()| — contiguous() 检查 + tensor 提取 \\
$\downarrow$ \\[4pt]
\textbf{Launch Wrapper} (\verb|softmax_kernel.cu|) \\
\\verb|switch(D)| 模板分发 + \\verb|<<<rows, 32>>>| 启动配置 \\\
$\downarrow$ \\[4pt]
\textbf{CUDA Kernel} (\verb|softmax_kernel.cu|, nvcc 编译) \\
\verb|softmax_warp_kernel<D>()| — Warp Shuffle 并行归约 \\
$\downarrow$ \\[4pt]
\textbf{GPU 硬件} (RTX 4060 Laptop, SM 8.9, 24 SMs, 8GB)
\end{tabular}
\end{center}

% ============================================================
\section{实现代码}
% ============================================================

\subsection{Naive 内核（串行基准）}

\textbf{设计思路：} 每个线程独立处理一行数据。Grid 维度为 \texttt{[ceil(rows/256)]}，Block 维度为 \texttt{[256]}。每线程串行遍历 $D$ 个元素三次：找最大值 $\to$ 算 $\exp$ 并求和 $\to$ 归一化。无线程间通信。

\begin{lstlisting}[style=cuda, caption={Naive Softmax 内核}, label=lst:naive]
template <int D>
__global__ void softmax_naive_kernel(
    const float* __restrict__ input,
    float* __restrict__ output,
    int rows, int stride)
{
    int row = blockIdx.x * blockDim.x + threadIdx.x;
    if (row >= rows) return;

    const float* x = input + row * stride;
    float* y = output + row * stride;

    // 第一遍: 找最大值 (Safe Softmax)
    float max_val = x[0];
    for (int i = 1; i < D; i++)
        max_val = fmaxf(max_val, x[i]);

    // 第二遍: 算 exp(x - max) 并求和
    float sum_val = 0.0f;
    for (int i = 0; i < D; i++) {
        y[i] = __expf(x[i] - max_val);
        sum_val += y[i];
    }

    // 第三遍: 归一化
    float inv_sum = 1.0f / sum_val;
    for (int i = 0; i < D; i++)
        y[i] *= inv_sum;
}
\end{lstlisting}

\textbf{时间复杂度：} 每线程 $O(3D)$。$D=1024$ 时，每线程 3072 次操作，完全串行。

\subsection{Warp 归约内核（并行优化）}

\textbf{设计思路：} 采用 \textbf{One-Warp-Per-Row} 映射——Grid 维度为 \texttt{[rows]}，Block 维度为 \texttt{[32]}（恰好一个 Warp）。32 个线程通过跨步访问分治处理 $D/32$ 个元素，然后通过 \texttt{\_\_shfl\_down\_sync} 蝴蝶归约和 \texttt{\_\_shfl\_sync} 广播完成全局归约。

\begin{lstlisting}[style=cuda, caption={Warp Shuffle Softmax 内核}, label=lst:warp]
template <int D>
__global__ void softmax_warp_kernel(
    const float* __restrict__ input,
    float* __restrict__ output,
    int rows, int stride)
{
    int row = blockIdx.x;               // 一行一个 Block
    if (row >= rows) return;

    const float* x = input + row * stride;
    float* y = output + row * stride;
    int lane = threadIdx.x;             // Warp 内编号 (0~31)

    // ===== Step 1: 跨步访问找局部最大值 =====
    // 数据分配: lane 0->x[0,32,64,...], lane 1->x[1,33,65,...]
    // 相邻线程访问相邻地址 -> 全局内存合并访问
    float max_val = -INFINITY;
    for (int i = lane; i < D; i += 32)
        max_val = fmaxf(max_val, x[i]);

    // ===== Step 2: Max Reduction + Broadcast =====
    // 蝴蝶归约: offset=16->8->4->2->1, 5 步合并 32 个局部值
    for (int offset = 16; offset > 0; offset /= 2)
        max_val = fmaxf(max_val,
            __shfl_down_sync(0xFFFFFFFF, max_val, offset));
    // Broadcast: 将 lane 0 手里的全局最大值复制给全部 32 线程
    max_val = __shfl_sync(0xFFFFFFFF, max_val, 0);

    // ===== Step 3: Sum Reduction + Broadcast =====
    float sum_val = 0.0f;
    for (int i = lane; i < D; i += 32) {
        float e = __expf(x[i] - max_val);
        y[i] = e;  sum_val += e;
    }
    for (int offset = 16; offset > 0; offset /= 2)
        sum_val += __shfl_down_sync(0xFFFFFFFF, sum_val, offset);
    sum_val = __shfl_sync(0xFFFFFFFF, sum_val, 0);

    // ===== Step 4: 归一化输出 =====
    float inv_sum = 1.0f / sum_val;
    for (int i = lane; i < D; i += 32)
        y[i] *= inv_sum;
}
\end{lstlisting}

\subsubsection{并行策略}

\begin{verbatim}
Warp 0 (32线程) -> 处理第 0 行
Warp 1 (32线程) -> 处理第 1 行
...
Warp M (32线程) -> 处理第 M 行

单个 Warp 内（跨步访问，合并访存）:
  lane  0 -> x[ 0], x[32], x[64], ...
  lane  1 -> x[ 1], x[33], x[65], ...
  ...
  lane 31 -> x[31], x[63], x[95], ...
  每线程只处理 D/32 个元素
\end{verbatim}

\subsubsection{蝴蝶归约过程}

\begin{verbatim}
初始化:     32 个线程各自有局部 max/sum
offset=16:  线程0和16合并, 1和17合并...  -> 剩 16 个
offset=8:   线程0和8合并,  1和9合并...   -> 剩 8 个
offset=4:   线程0和4合并,  1和5合并...   -> 剩 4 个
offset=2:   线程0和2合并,  1和3合并     -> 剩 2 个
offset=1:   线程0和1合并                -> 剩 1 个 [在 lane 0]
Broadcast:  lane 0 -> 全部 32 线程
\end{verbatim}

\textbf{关键机制：} 蝴蝶归约后结果仅保存在 lane 0，必须通过 \texttt{\_\_shfl\_sync} 广播给 Warp 内全部 32 条线程，所有线程才能使用正确的全局值进行后续归一化。

\textbf{时间复杂度：} 每线程 $O(3D/32)$ 次内存访问 + $5 \times 2$ 次 shuffle 通信。

\subsection{Launch 包装器（模板分发）}

\begin{lstlisting}[style=cuda, caption={Launch Wrapper}]
void softmax_warp_launch(const float* input, float* output,
                         int rows, int D, int stride) {
    // Grid: [rows], Block: [32] — 一行恰好一个 Warp
    switch (D) {
        case 64:  softmax_warp_kernel<64>
                      (*@<<<@*)rows, 32(*@>>>@*)(input, output, rows, stride); break;
        case 128: softmax_warp_kernel<128>
                      (*@<<<@*)rows, 32(*@>>>@*)(input, output, rows, stride); break;
        case 256: softmax_warp_kernel<256>
                      (*@<<<@*)rows, 32(*@>>>@*)(input, output, rows, stride); break;
        case 512: softmax_warp_kernel<512>
                      (*@<<<@*)rows, 32(*@>>>@*)(input, output, rows, stride); break;
        case 1024:softmax_warp_kernel<1024>
                      (*@<<<@*)rows, 32(*@>>>@*)(input, output, rows, stride); break;
        default:  fprintf(stderr, "Unsupported D=%d\n", D); exit(1);
    }
    CUDA_CHECK(cudaGetLastError());
    CUDA_CHECK(cudaDeviceSynchronize());
}
\end{lstlisting}

\subsection{PyTorch 绑定}

\textbf{C++ 绑定} (\texttt{csrc/softmax\_bindings.cpp})：

\begin{lstlisting}[style=cuda, caption={C++ Torch 绑定}]
torch::Tensor softmax_forward_warp(torch::Tensor input) {
    input = input.contiguous();
    auto sizes = input.sizes();
    int D = sizes.back();
    int rows = 1;
    for (size_t i = 0; i < sizes.size() - 1; i++) rows *= sizes[i];
    auto output = torch::empty_like(input);
    softmax_warp_launch(input.data_ptr<float>(), output.data_ptr<float>(),
                        rows, D, D);
    return output;
}

PYBIND11_MODULE(TORCH_EXTENSION_NAME, m) {
    m.def("softmax_forward_naive", &softmax_forward_naive);
    m.def("softmax_forward_warp", &softmax_forward_warp);
}
\end{lstlisting}

\textbf{Python API} (\texttt{softmax/\_\_init\_\_.py})：

\begin{lstlisting}[style=python, caption={Python API}]
from . import _C

def softmax_warp(x):
    """Warp-reduction CUDA softmax forward.
    One warp (32 threads) per row.
    """
    return _C.softmax_forward_warp(x)
\end{lstlisting}

\textbf{构建配置} (\texttt{setup.py})：

\begin{lstlisting}[style=python, caption={setup.py 编译配置}]
CUDAExtension(
    "softmax._C",
    ["csrc/softmax_bindings.cpp", "csrc/softmax_kernel.cu"],
    extra_compile_args={
        "cxx": ["-O3", "-fopenmp"],
        "nvcc": ["-O3", "--expt-relaxed-constexpr",
                 "-gencode=arch=compute_89,code=sm_89",
                 "-gencode=arch=compute_80,code=sm_80"],
    },
)
\end{lstlisting}

% ============================================================
\section{实验设置}
% ============================================================

\subsection{测试环境}

\begin{table}[H]
\centering
\caption{测试环境配置}
\begin{tabular}{ll}
\toprule
\textbf{项目} & \textbf{配置} \\
\midrule
GPU & NVIDIA GeForce RTX 4060 Laptop GPU \\
计算能力 & SM 8.9 (Ada Lovelace) \\
SM 数量 & 24 \\
显存 & 8.6 GB \\
CUDA Toolkit & 13.3 \\
PyTorch & 2.11.0+cu130 \\
Host 编译器 & g++ 15.2 (conda) \\
数据类型 & float32 \\
\bottomrule
\end{tabular}
\end{table}

\subsection{测试方法}

\begin{itemize}[itemsep=3pt]
    \item \textbf{对比对象：} Naive（手写串行）、Warp（手写并行归约）、torch.softmax（PyTorch 官方，底层 cuDNN）
    \item \textbf{测试维度：} $D \in \{64, 128, 256, 512, 1024\}$
    \item \textbf{数据规模：} $D \leq 512$ 时 $\text{rows} = 32768$，$D = 1024$ 时 $\text{rows} = 16384$
    \item \textbf{预热：} 5 轮
    \item \textbf{取平均：} 30 轮
    \item \textbf{正确性验证：} 以 \texttt{torch.softmax} 为参考，计算最大绝对偏差
\end{itemize}

\subsection{测试脚本}

\begin{lstlisting}[style=python, caption={性能测试脚本 (benchmark.py)}]
import torch, time
from softmax import softmax_naive, softmax_warp

WARMUP, ITERS = 5, 30

def benchmark(dims):
    for D in dims:
        rows = 32768 if D <= 512 else 16384
        x = torch.randn(rows, D, device='cuda', dtype=torch.float32)

        # Warmup
        for _ in range(WARMUP):
            softmax_naive(x); softmax_warp(x); torch.softmax(x, dim=-1)
        torch.cuda.synchronize()

        # Naive
        t0 = time.perf_counter()
        for _ in range(ITERS): softmax_naive(x)
        torch.cuda.synchronize()
        t_naive = (time.perf_counter() - t0) / ITERS * 1e6

        # Warp
        t0 = time.perf_counter()
        for _ in range(ITERS): softmax_warp(x)
        torch.cuda.synchronize()
        t_warp = (time.perf_counter() - t0) / ITERS * 1e6

        # torch
        t0 = time.perf_counter()
        for _ in range(ITERS): torch.softmax(x, dim=-1)
        torch.cuda.synchronize()
        t_torch = (time.perf_counter() - t0) / ITERS * 1e6

        # Correctness
        ref = torch.softmax(x, dim=-1)
        d_naive = (softmax_naive(x) - ref).abs().max().item()
        d_warp  = (softmax_warp(x) - ref).abs().max().item()

        yield {'D':D, 'naive_us':t_naive, 'warp_us':t_warp,
               'torch_us':t_torch, 'naive_err':d_naive, 'warp_err':d_warp}
\end{lstlisting}

编译与运行：
\begin{lstlisting}[style=shell]
source env.sh
python3 setup.py build_ext --inplace
python3 softmax/benchmark/benchmark.py
\end{lstlisting}

% ============================================================
\section{实验结果}
% ============================================================

\subsection{正确性验证}

以 \texttt{torch.softmax} 结果为参考，计算最大绝对偏差：

\begin{table}[H]
\centering
\caption{正确性验证结果}
\begin{tabular}{cccc}
\toprule
$D$ & Naive 偏差 & Warp 偏差 & 结论 \\
\midrule
64  & $1.49\times10^{-7}$ & $5.96\times10^{-8}$ & $\checkmark$ \\
128 & $1.49\times10^{-7}$ & $5.96\times10^{-8}$ & $\checkmark$ \\
256 & $1.04\times10^{-7}$ & $2.98\times10^{-8}$ & $\checkmark$ \\
512 & $1.49\times10^{-7}$ & $1.49\times10^{-8}$ & $\checkmark$ \\
1024& $9.31\times10^{-8}$ & $7.45\times10^{-9}$ & $\checkmark$ \\
\bottomrule
\end{tabular}
\end{table}

所有维度下最大绝对偏差 $<2\times10^{-7}$，远优于 $10^{-5}$ 容限。Warp 版本的精度（$\sim5\times10^{-8}$）系统性地优于 Naive 版本（$\sim1.5\times10^{-7}$），因为归约将浮点累加次数从 $O(D)$ 降至 $O(D/32 + \log_2 32)$，减少了舍入误差累积。

\subsection{延迟对比}

处理 $\text{rows} \times D$ 矩阵的平均耗时（$\mu\text{s}$，越小越好）：

\begin{table}[H]
\centering
\caption{延迟对比}
\begin{tabular}{cccccc}
\toprule
$D$ & Naive ($\mu$s) & Warp ($\mu$s) & torch ($\mu$s) & Warp/Naive & Warp/torch \\
\midrule
64   & 520  & 131  & 32   & 4.0$\times$ & 4.07$\times$ \\
128  & 970  & 129  & 48   & 7.5$\times$ & 2.68$\times$ \\
256  & 1840 & 313  & 283  & 5.9$\times$ & 1.11$\times$ \\
512  & 3185 & 663  & 590  & 4.8$\times$ & 1.12$\times$ \\
1024 & 5651 & 1319 & 1172 & 4.3$\times$ & 1.13$\times$ \\
\bottomrule
\end{tabular}
\end{table}

\subsection{吞吐量对比}

每秒可处理的行数（M rows/s，越大越好）：

\begin{table}[H]
\centering
\caption{吞吐量对比}
\begin{tabular}{cccc}
\toprule
$D$ & Naive (M rows/s) & Warp (M rows/s) & torch (M rows/s) \\
\midrule
64   & 63.03  & 250.13 & 1019.04 \\
128  & 33.77  & 254.00 & 679.88  \\
256  & 17.81  & 104.66 & 115.80  \\
512  & 10.29  & 49.43  & 55.56   \\
1024 & 5.80   & 24.84  & 27.97   \\
\bottomrule
\end{tabular}
\end{table}

% ============================================================
\section{分析与讨论}
% ============================================================

\subsection{Warp vs Naive}

Warp Shuffle 归约在所有测试维度上均显著优于 Naive 串行实现，加速比 $4$--$8\times$。加速来源于两点：

\begin{enumerate}[itemsep=4pt]
    \item \textbf{分治减少单线程工作量：} Naive 每线程处理 $D$ 个元素，Warp 版每线程仅处理 $D/32$ 个元素
    \item \textbf{Warp Shuffle 是硬件级通信：} \texttt{\_\_shfl\_down\_sync} 直接通过片上连线交换寄存器数据，延迟约 1 个时钟周期。5 步归约总计约 5 cycles，而全局内存访问需要 $\sim$400 cycles
\end{enumerate}

\subsection{Warp vs torch.softmax（CUDA cuDNN）}

\begin{table}[H]
\centering
\caption{Warp vs torch 详细对比}
\begin{tabular}{cccc}
\toprule
$D$ & 手写 Warp ($\mu$s) & torch/cuDNN ($\mu$s) & 差距 \\
\midrule
256  & 313  & 283  & $1.11\times$ \\
512  & 663  & 590  & $1.12\times$ \\
1024 & 1319 & 1172 & $1.13\times$ \\
\bottomrule
\end{tabular}
\end{table}

$D \geq 256$ 时手写内核与 CUDA 自带函数性能差距 $<15\%$，$D=1024$ 时仅差 $13\%$。这表明：

\begin{itemize}[itemsep=3pt]
    \item 手写 Warp Shuffle 归约已接近 GPU 硬件极限
    \item cuDNN 的 softmax 内核与我们采用同一种算法模式（warp reduction），主要区别在于 cuDNN 使用了 \texttt{float4} 向量化加载（每次读 4 个 float）和针对特定 $D$ 值的汇编级优化
    \item 但 $D \geq 256$ 时，计算瓶颈从内存带宽转向计算吞吐，向量化优势递减
\end{itemize}

\subsection{$D$ 较小时 torch 的优势}

$D=64$ 时 torch 仍有 $4.07\times$ 优势。因为该维度下每线程的工作量（$64/32=2$ 个元素）极低，kernel 启动开销和 warp 调度开销相对占比大。torch 的 kernel 针对小 $D$ 有专门优化（如融合相邻行、使用更小的 block）。

这也是手写 kernel 的下一步优化方向：引入 \texttt{float4} 向量化全局加载。

\subsection{与 CUB/Cooperative Groups 的关系}

手写的 Reduce + Broadcast 模式在 CUDA 生态中有对应的库实现：

\begin{lstlisting}[style=cuda]
// CUB (CUDA UnBound)
cub::BlockReduce<float>(smem).Reduce(val, cub::Sum());

// Cooperative Groups
auto tile = cg::tiled_partition<32>(cg::this_thread_block());
cg::reduce(tile, val, cg::plus<float>());
\end{lstlisting}

三者底层均使用相同的 \texttt{\_\_shfl\_down\_sync} 硬件指令。手写版本的优势在于零依赖、完全透明、灵活控制归约策略。

% ============================================================
\section{结论}
% ============================================================

\begin{enumerate}[itemsep=6pt]
    \item \textbf{Warp Shuffle 是 GPU 上最快的线程间通信机制。} 利用蝴蝶归约 + 广播模式，手写的 Softmax 内核在 $D \geq 256$ 时与 CUDA 自带的 cuDNN 实现性能持平（差距 $<15\%$）。

    \item \textbf{80 行 CUDA 代码达到接近硬件极限的性能。} $D=1024$ 时，手写 Warp 内核仅比工业级 cuDNN 慢 $13\%$，充分证明了 Warp Shuffle 归约模式的高效性。

    \item \textbf{归约减少浮点误差。} Warp 版本的数值精度系统性地优于 Naive 版本，因为归约将累加次数从 $O(D)$ 降至 $O(D/32 + \log_2 32)$。

    \item \textbf{$D$ 较小时的性能差距来自向量化加载。} 引入 \texttt{float4} 后预期可缩小 $D=64/128$ 时的差距至 $2\times$ 以内。

    \item \textbf{Warp Shuffle 归约模式具有普适性。} 不限于 Softmax——任何跨线程聚合操作（求和、最大值、最小值、前缀和、LayerNorm、RMSNorm）都可以用相同的 Reduce + Broadcast 模式高效实现。
\end{enumerate}

% ============================================================
\appendix
\section{项目文件结构}
% ============================================================

\begin{verbatim}
softmax/
├── __init__.py                     # Python API
├── benchmark/
│   ├── benchmark.py                # 性能测试脚本
│   └── report.md / report.tex      # 本报告
csrc/
├── softmax_kernel.cu               # CUDA 内核实现 (Naive + Warp)
└── softmax_bindings.cpp            # PyTorch C++ 绑定
setup.py                            # 编译配置
docs/softmax/
├── ppt/index.html                  # 课堂演示网页 PPT
├── ppt/script.md                   # 讲稿
└── implementation.md               # 详细实现文档
\end{verbatim}

\section{完整运行输出}

\begin{lstlisting}[style=shell, caption={benchmark.py 完整输出}]
======================================================================
Softmax Performance Benchmark
======================================================================
GPU: NVIDIA GeForce RTX 4060 Laptop GPU
PyTorch: 2.11.0+cu130, CUDA: 13.0

     D | Naive (us) |  Warp (us) | torch (us) | Warp/Naive | Warp/torch
----------------------------------------------------------------------
    64 |        520 |        131 |         32 |       4.0x |      4.07x
   128 |        970 |        129 |         48 |       7.5x |      2.68x
   256 |       1840 |        313 |        283 |       5.9x |      1.11x
   512 |       3185 |        663 |        590 |       4.8x |      1.12x
  1024 |       5651 |       1319 |       1172 |       4.3x |      1.13x

     D |    Naive err |     Warp err |   Status
--------------------------------------------------
    64 |     1.49e-07 |     5.96e-08 |        PASS
   128 |     1.49e-07 |     5.96e-08 |        PASS
   256 |     1.04e-07 |     2.98e-08 |        PASS
   512 |     1.49e-07 |     1.49e-08 |        PASS
  1024 |     9.31e-08 |     7.45e-09 |        PASS

     D | Naive (M rows/s) |  Warp (M rows/s) | torch (M rows/s)
------------------------------------------------------------
    64 |            63.03 |           250.13 |          1019.04
   128 |            33.77 |           254.00 |           679.88
   256 |            17.81 |           104.66 |           115.80
   512 |            10.29 |            49.43 |            55.56
  1024 |             5.80 |            24.84 |            27.97

=== Environment ===
GPU: NVIDIA GeForce RTX 4060 Laptop GPU
Compute Capability: (8, 9)
VRAM: 8.6 GB
CUDA: 13.0
PyTorch: 2.11.0+cu130
Warmup: 5, Iters: 30, dtype: float32
\end{lstlisting}

\end{document}
